\section{Conclusion}

\subsection{Main Conclusion and Hypothesis Evaluation}
This research critically evaluated the relative accuracy of vision-based and GPS-based speed estimation algorithms for autonomous vehicles. The principal inquiry examined the performance of a vision-based model, employing Region of Interest (ROI) Dense Optical Flow and dynamic affine calibration, in comparison to conventional GPS probe data against a CAN-bus ground truth.

The principal conclusion is unequivocal: consumer-grade GPS trajectory data markedly surpasses the proposed pixel-tracking vision methodology in terms of both accuracy and stability. The GPS-based pipeline attained a Mean Absolute Error (MAE) of 1.68 km/h, whereas the vision-based system yielded an MAE of 4.50 km/h, demonstrating that the optical flow approach is nearly three times less precise.

These findings definitively refute the original research hypothesis, which proposed that the vision-based method would result in a smaller error margin owing to satellite latency. The data indicate that, although GPS exhibits a slight negative bias of -1.43 km/h, its overall ability to accurately ascertain a vehicle’s kinematics in typical urban environments is significantly superior to that of Dense Optical Flow.

\subsection{Shortcomings in the Methodology}
While the setup benefited from the highly synchronised ground truth provided by the \textit{comma2k19} dataset \cite{schafer2018commute}, a thorough evaluation reveals several important methodological limitations.

\subsubsection{Algorithmic Fragility}
Relying on the Farnebäck algorithm meant tracking raw visual shifts rather than physical objects. The methodology failed to account for environmental visual noise, such as bridge shadows and oncoming traffic. This lack of semantic understanding caused the algorithm to fail at detecting deceleration, resulting in a systemic overestimation of speed.

\subsubsection{Limited Generalisability}
The dataset exhibited well-illuminated, clear Californian weather conditions. As the optical flow methodology deteriorates during slight variations in lighting even under optimal circumstances, its efficacy in adverse weather conditions or during nocturnal driving would be substantially compromised.

\subsubsection{Camera Calibration Assumptions}
The dynamic affine calibration mathematically presumes a stable camera pitch. In practice, intense braking induces suspension dive, which alters the pitch and compromises the integrity of the optical flow vectors. The methodology did not incorporate pitch-compensation algorithms to address this physical variable.

\subsection{Suggestions for Further Research}
To further the development of Intelligent Transportation Systems (ITS), several promising avenues for additional research present themselves:

\subsubsection{Transition to Deep Learning (YOLO)}
Future vision-based research should incorporate deep learning frameworks such as YOLO (You Only Look Once) for object detection. By precisely monitoring the physical boundaries of a vehicle, the system attains a semantic understanding, thereby enabling it to disregard peripheral visual noise and shadows that previously compromised the effectiveness of optical flow techniques.

\subsubsection{Sensor Fusion via Kalman Filtering}
Research should shift from assessing GPS and vision as mutually exclusive solutions towards their active integration. GPS offers reliable positioning but is subject to slight latency, whereas camera data provides immediate responsiveness but remains environmentally sensitive. The utilisation of a Kalman Filter to combine these data streams could theoretically attain error margins below 1 km/h.

\subsubsection{Algorithmic Noise Mitigation for GPS}
Considering that raw GPS data demonstrated exceptional performance, future research efforts should focus on employing Deep Neural Networks (DNNs) with standard satellite trajectories in order to predict and mitigate multipath signal reflections, thereby advancing consumer-grade receivers towards near-Real Time Kinematic (RTK) accuracy.